% =====================================================================
% Công trình liên quan. Bài mới trích nằm trong refs-related.bib (đã đối chiếu tên, tác giả, nơi đăng).
% Mục này không trích con số nào từ các bài báo.
\section{Công trình liên quan}
\label{sec:related}

\subsection{Truy vấn văn bản--ảnh và giới hạn của CLIP}

Các mô hình so khớp ảnh--văn bản trước CLIP căn từng vùng ảnh với từng từ trong câu; ví dụ, SCAN~\cite{rw-lee2018scan}
dùng cross-attention giữa vùng ảnh và từ để tìm các cặp tương ứng. CLIP~\cite{radford2021clip} và
OpenCLIP~\cite{cherti2023openclip} huấn luyện tương phản (contrastive) hai bộ mã hoá ảnh và văn bản trên các cặp ảnh--câu
quy mô lớn. Truy vấn khi đó là tìm láng giềng gần theo cosine, làm được nhanh với FAISS~\cite{johnson2019faiss}, nhưng
mỗi bên chỉ còn một vector. ARO~\cite{rw-yuksekgonul2023aro} cho thấy các mô hình này đạt kết quả truy vấn tốt trên các
benchmark sẵn có mà không cần dùng thông tin về thành phần và thứ tự trong câu, và kém ở quan hệ, ở gắn thuộc tính với vật
thể và ở thứ tự; bài báo sửa bằng mẫu âm khó khi huấn luyện. Winoground~\cite{rw-thrush2022winoground} kiểm tra cùng điểm
yếu bằng các cặp câu có cùng các từ nhưng khác thứ tự. Nhóm không huấn luyện lại CLIP mà đưa cấu trúc vào ở bước xếp lại.

\subsection{Scene graph và GNN trong truy vấn ảnh}

Johnson và cộng sự~\cite{rw-johnson2015sgretrieval} đề xuất truy vấn ảnh bằng scene graph: truy vấn là một đồ thị gồm vật
thể, thuộc tính và quan hệ; một trường ngẫu nhiên có điều kiện (CRF) tìm cách gắn đồ thị vào từng ảnh, và độ khả dĩ của
cách gắn đó là điểm xếp hạng. Visual Genome~\cite{krishna2017vg} mở rộng ý tưởng này thành một bộ ảnh lớn có scene graph
gán tay. Các mô hình so khớp sau đó đưa cấu trúc vào bộ mã hoá. VSRN~\cite{rw-li2019vsrn} nối các vùng ảnh phát hiện được
thành đồ thị rồi chạy GCN. SGM~\cite{rw-wang2020sgm} biểu diễn cả ảnh và câu bằng scene graph và so khớp ở mức vật thể và
mức quan hệ. GSMN~\cite{rw-liu2020gsmn} so khớp ở mức node và mức cấu trúc. Structure-CLIP~\cite{rw-huang2024structureclip}
dùng scene graph của câu để sinh mẫu âm khó và đưa tri thức scene graph vào khi tinh chỉnh CLIP. Phía câu, scene graph được
tách bằng parser dựa trên luật và cây cú pháp~\cite{schuster2015sgparser} hoặc bằng parser seq2seq như
FACTUAL~\cite{li2023factual}. Các lớp GNN thường dùng là GCN~\cite{rw-kipf2017gcn}, GraphSAGE~\cite{rw-hamilton2017sage},
GAT~\cite{rw-velickovic2018gat} và GATv2~\cite{brody2022gatv2}. Khác với các mô hình trên, nhóm giữ nguyên hai bộ mã hoá
của CLIP; GATv2 chỉ chấm lại 50 ứng viên, và điểm đồ thị tách được thành từng cặp phần khớp giữa câu và ảnh, nên lời giải
thích là chính các số hạng của điểm.

\subsection{Đồ thị tri thức đa phương thức}

Các mô hình embedding đồ thị tri thức học vector cho thực thể và quan hệ để chấm điểm một bộ ba:
TransE~\cite{rw-bordes2013transe} coi quan hệ là một phép tịnh tiến, RotatE~\cite{sun2019rotate} coi quan hệ là một phép
quay trong không gian phức. Đồ thị tri thức đa phương thức (multimodal knowledge graph, MMKG) gắn thêm ảnh và mô tả cho
thực thể~\cite{rw-chen2024mmkgsurvey}. IKRL~\cite{rw-xie2017ikrl} học biểu diễn thực thể từ cả triple lẫn ảnh. Các mô hình
hoàn thiện MMKG (MMKGC) gần đây gồm MKGformer~\cite{rw-chen2022mkgformer} (transformer lai, trộn hai phương thức ở nhiều
mức), MMRNS~\cite{xu2022mmrns} (lấy mẫu âm theo quan hệ; đây cũng là nguồn của MKG-W), NativE~\cite{zhang2024native} và
AdaMF-MAT~\cite{rw-zhang2024adamf} (xử lý thực thể thiếu hoặc mất cân bằng phương thức), và mô hình trên không gian đa độ
cong (mixed-curvature) của~\cite{rw-gao2025mixedcurvature}. Các mô hình này được đánh giá bằng link prediction, không phải
bằng truy vấn ảnh. Nhóm tái lập NativE, rồi dùng embedding và hàm chấm bộ ba của nó làm tín hiệu xếp lại, đặt cạnh một tín
hiệu đơn giản hơn: số cạnh tri thức giữa thực thể ứng viên và các thực thể được nhắc trong câu.

\subsection{Giải thích GNN và đánh giá lời giải thích}

Phần lớn các cách giải thích GNN là hậu kiểm (post-hoc): sau khi mô hình đã dự đoán, chúng tìm phần đồ thị quan trọng bằng
mask học được (GNNExplainer~\cite{xai-ying2019gnnexplainer}, PGExplainer~\cite{xai-luo2020pgexplainer},
GraphMask~\cite{xai-schlichtkrull2021graphmask}), bằng giá trị Shapley (GraphSVX~\cite{xai-duval2021graphsvx},
SubgraphX~\cite{xai-yuan2021subgraphx}, GStarX~\cite{xai-zhang2022gstarx}), bằng lan truyền relevance
(GNN-LRP~\cite{xai-schnake2022gnnlrp}) hoặc bằng phản thực tế (CF-GNNExplainer~\cite{xai-lucic2022cfgnn}). Độ trung thực
(fidelity) thường đo bằng cách bỏ, hoặc chỉ giữ, phần được giải thích rồi xem dự đoán thay đổi bao
nhiêu~\cite{xai-yuan2023survey}. Fid$_\alpha$~\cite{xai-zheng2024fidalpha} chỉ bỏ ngẫu nhiên một phần nhỏ để đồ thị sau khi
sửa không lệch xa phân phối lúc huấn luyện; GraphXAI~\cite{xai-agarwal2023graphxai} chấm lời giải thích bằng đáp án biết
trước. Jain và Wallace~\cite{xai-jain2019attention} lưu ý rằng trọng số attention không nhất thiết là lời giải thích. Các
cách trên chủ yếu viết cho phân loại node hoặc đồ thị; RankingSHAP~\cite{xai-heuss2025rankingshap} đưa cách tính kiểu
Shapley sang mô hình xếp hạng. Vì đề bài không chấp nhận lời giải thích dựng sau khi đã có kết quả, lời giải thích chính
của nhóm là phép phân rã của chính điểm xếp hạng; các cách hậu kiểm chỉ dùng làm mốc so sánh trên cùng độ đo.

\subsection{Encoder văn bản đa ngữ}

CLIP ViT-B/32 của OpenAI chỉ học văn bản tiếng Anh. Có hai cách thường dùng để truy vấn bằng ngôn ngữ khác: thay tháp văn
bản bằng XLM-R~\cite{vn-conneau2020xlmr} rồi huấn luyện tương phản với tháp ảnh trên LAION-5B~\cite{vn-schuhmann2022laion5b}
(bản OpenCLIP~\cite{vn-ilharco2021openclip}), hoặc chưng cất (knowledge distillation) một bộ mã hoá văn bản đa ngữ vào
không gian của CLIP~\cite{vn-reimers2020multilingual}. Nhóm so cả hai cách trên MKG-W val trước khi chạy trên MMKG tiếng
Việt (Mục~\ref{sec:vn-encoder}).
